% ===========
%  CHAPTER 3
% ===========
\OpenGrid

\ParallelChapter
  {Reasoning and Sense Making in the Curriculum}
  {Luận và hiểu trong chương trình học}
  {}

\TriPara
  {
    \EN{\emph{Reasoning and sense making are integral to the experiences of all students across all areas of the high school mathematics curriculum.}}
  }
  {
    \VI{\emph{Luận và hiểu là phần không thể tách rời trong trải nghiệm của mọi học sinh, xuyên suốt mọi mảng của chương trình toán học trung học phổ thông.}}
  }{\NOTE{}}

\TriPara
  {
    \EN{Reasoning and sense making should be pervasive across all areas of the high school mathe matics curriculum. Although aspects of formal reasoning are frequently emphasized in geometry, students are less likely to encounter reasoning in other areas of the curriculum, such as algebra. When reasoning and sense making are infused throughout the curriculum, they lend coherence across the domains of mathematics---number, algebra, geometry, and statistics-whether the curriculum is arranged in the customary U.S. course---based sequence (first-year algebra, geometry, second-year algebra) or in an integrated manner. A focus on reasoning also helps students see how new concepts connect with existing knowledge.}
  }
  {
    \VI{Luận và hiểu cần hiện diện xuyên suốt mọi mảng của chương trình toán học trung học phổ thông. Mặc dù những phương diện của luận hình thức thường được nhấn mạnh trong hình học, học sinh lại ít có cơ hội gặp luận trong những mảng khác của chương trình, chẳng hạn như đại số. Khi luận và hiểu được thấm vào toàn bộ chương trình học, chúng tạo nên sự nhất quán giữa các lĩnh vực của toán học---số, đại số, hình học và thống kê---bất kể chương trình được sắp xếp theo trình tự khóa học quen thuộc ở Hoa Kì (đại số năm thứ nhất, hình học, đại số năm thứ hai) hay theo cách tích hợp. Việc chú trọng vào luận cũng giúp học sinh thấy các khái niệm mới kết nối với kiến thức hiện có như thế nào.}
  }{\NOTE{}}

\TriPara
  {
    \EN{We emphasize that we do not propose introducing ``reasoning'' as a set of topics to be added to a crowded curriculum, but as a stance toward learning mathematics. Developing strong reasoning habits will of course take instructional time. However, it also promises compensating efficiencies. First, the initial part of each course that is frequently spent reteaching content from previous courses may be reduced if those courses emphasize reasoned connections with existing knowledge, so that students are better able to retain what they have learned. Second, emphasizing the underlying connections that promote reasoning and sense making introduces coherence that allows streamlining of the curriculum. As attention turns to how those underlying connections build across the high school years, less time needs to be spent addressing lists of particular skills that need to be mastered.}
  }
  {
    \VI{Chúng tôi nhấn mạnh rằng chúng tôi không đề xuất đưa ``luận'' vào như một tập hợp chủ đề mới bổ sung cho chương trình vốn đã quá tải, mà như một lập trường đối với việc học toán. Phát triển thói quen luận vững chắc dĩ nhiên sẽ đòi hỏi thời gian dạy học. Tuy nhiên, điều đó cũng hứa hẹn hiệu quả bù lại. Thứ nhất, phần đầu mỗi khóa học, vốn thường được dùng để dạy lại nội dung từ khóa học trước, có thể được rút ngắn nếu các khóa học đó nhấn mạnh những kết nối có cơ sở luận với kiến thức hiện có, nhờ đó học sinh có khả năng ghi nhớ tốt hơn những gì đã học. Thứ hai, việc nhấn mạnh các kết nối nền tảng thúc đẩy luận và hiểu sẽ tạo ra tính nhất quán, qua đó cho phép tinh giản chương trình. Khi sự chú ý chuyển sang cách những kết nối nền tảng ấy được xây dựng dần qua các năm trung học phổ thông, sẽ cần dành ít thời gian hơn cho việc xử lí những danh sách kĩ năng cụ thể mà học sinh cần nắm vững.}
  }
  {
    \NOTE{Điều cần được chặn lại ngay từ đầu là cách hiểu rằng muốn tăng cường luận thì phải thêm một mảng nội dung mới vào chương trình đã quá tải. Tác giả bác bỏ trực tiếp cách hiểu ấy. Luận không được đặt ra như một ``phần thêm'', mà như một lập trường đối với việc học toán: học bằng cách kết nối cái mới với cái đã biết, thấy các mối quan hệ đan xen hỗ trợ nhau thay vì những mảnh vụn rời rạc. Vì thế, thời gian dành cho việc phát triển thói quen luận không phải là phần thời gian bị lấy mất khỏi chương trình, mà là khoản đầu tư làm đổi cấu trúc vận hành của toàn bộ việc dạy học.

    Luận có thể tốn thời gian trước mắt, nhưng lại tiết kiệm thời gian về sau theo hai hướng. Một là giảm nhu cầu dạy lại ở đầu mỗi khóa học, vì kiến thức đã học trước đó không còn được giữ bằng trí nhớ thủ tục đơn lẻ mà bằng các kết nối có nghĩa, nên bền hơn và dễ gọi lại hơn. Hai là tạo ra tính nhất quán cho chương trình, nhờ đó trọng tâm dần chuyển từ việc xử lí những danh sách kĩ năng rời rạc sang việc xây dựng các kết nối nền tảng xuyên suốt những năm THPT. Ở đây, các tác giả không chỉ bảo vệ giá trị nhận thức của luận và hiểu, mà còn bảo vệ chúng bằng lí do rất thực tiễn: chúng làm cho chương trình gọn hơn, chặt hơn và bớt lãng phí hơn.}
  }

\TriPara
  {
    \EN{For example, teaching students to factor polynomials often consumes considerable time, as students are taught methods for factoring (1) common monomial terms from a polynomial, (2) trinomials with leading coefficient of 1 (with or without negative coefficients), (3) trinomials with a leading coefficient other than 1 (which may eventually be a noninteger), (4) special cases, such as perfect squares and the difference of squares (or cubes), and (5) more-complex trinomials in which the exponent of the variable in one term is double that in another. Nonetheless, students who are given experiences that help them see all these different forms as results of using the distributive law may be more likely to understand the common structure behind the different factoring methods, and thus develop fluency with them. The area model is particularly useful in visualizing the distributive law; see examples 4 and 8 for a discussion of area models. This approach may also facilitate long-term retention more effectively than teaching particular procedures for factoring in isolation.}
  }
  {
    \VI{Ví dụ, việc dạy học sinh phân tích đa thức thành nhân tử thường tiêu tốn khá nhiều thời gian, vì các em được dạy những phương pháp phân tích: (1) đưa nhân tử đơn thức chung ra ngoài; (2) phân tích tam thức có hệ số đầu bằng 1 (có hoặc không có hệ số âm); (3) phân tích tam thức có hệ số đầu khác 1, mà về sau hệ số này thậm chí có thể không nguyên; (4) các trường hợp đặc biệt như bình phương hoàn chỉnh và hiệu của hai bình phương (hoặc hai lập phương); và (5) những tam thức phức tạp hơn, trong đó số mũ của biến ở một hạng tử gấp đôi số mũ của biến ở hạng tử khác. Tuy nhiên, nếu học sinh có những trải nghiệm giúp các em thấy rằng tất cả dạng khác nhau ấy đều là kết quả của việc vận dụng luật phân phối, thì các em có nhiều khả năng hiểu được cấu trúc chung nằm sau các phương pháp phân tích khác nhau, và từ đó phát triển sự thành thạo với chúng. Mô hình diện tích đặc biệt hữu ích trong việc trực quan hóa luật phân phối; xem Ví dụ 4 và 8 để biết thêm về mô hình diện tích. Cách tiếp cận này cũng có thể hỗ trợ việc ghi nhớ lâu dài hiệu quả hơn so với việc dạy tách biệt từng thủ tục phân tích thành nhân tử.}
  }
  {
    \NOTE{Phân tích thành nhân tử thường được dạy như một danh sách các dạng riêng: đặt nhân tử chung, tam thức bậc hai với hệ số đầu bằng $1$, khác $1$, hiệu hai bình phương, hay những dạng có thể quy về một ẩn phụ. Đoạn này muốn kéo các dạng đó về một cấu trúc chung: chúng đều có thể được hiểu như kết quả của việc đi ngược lại luật phân phối. Chẳng hạn, từ

    $$
    a(b+c)=ab+ac
    $$

    ta có chiều ngược lại

    $$
    ab+ac=a(b+c).
    $$

    Theo cách nhìn này, các thủ tục phân tích không còn là những mẹo rời rạc, mà là những biến thể của một cái chung.

    Điểm tác giả muốn nhấn mạnh là: khi học sinh thấy được cấu trúc chung, các em không chỉ làm đúng từng dạng bài, mà còn hiểu vì sao các cách phân tích khác nhau lại có liên hệ với nhau. Nhờ đó, sự thành thạo không chỉ dựa vào ghi nhớ thủ tục, mà dựa vào việc nhận ra cấu trúc. Mô hình diện tích được nhắc đến vì nó giúp trực quan hóa luật phân phối, nên đặc biệt hữu ích trong việc làm lộ cấu trúc chung ấy. Theo lập trường của đoạn văn, dạy theo hướng đó có thể hỗ trợ ghi nhớ lâu dài hiệu quả hơn so với việc dạy tách biệt từng thủ tục phân tích thành nhân tử.}
  }

\TriPara
  {
    \EN{The following chapters in this section illustrate how the high school mathematics curriculum can be focused on broad themes that promote reasoning and sense making within five specific content areas of the high school curriculum:

    \begin{itemize}
      \item Reasoning with Number and Measurement
      
      \item Reasoning with Algebraic Symbols
      
      \item Reasoning with Functions
      
      \item Reasoning with Geometry
      
      \item Reasoning with Statistics and Probability
    \end{itemize}}
  }
  {
    \VI{Các chương tiếp theo trong phần này minh họa cách chương trình toán học trung học phổ thông có thể được định hướng theo những chủ đề rộng thúc đẩy luận và hiểu trong năm mảng nội dung cụ thể của chương trình học trung học phổ thông:

    \begin{itemize}
      \item Luận với số và đo lường
      
      \item Luận với kí hiệu đại số
      
      \item Luận với hàm số
      
      \item Luận với hình học
      
      \item Luận với thống kê và xác suất
    \end{itemize}}
  }
  {
    \NOTE{Năm mảng được nêu ở đây không nên hiểu là năm chuẩn nội dung chính thức, cũng không phải là năm phân môn đầy đủ của toán học trung học phổ thông. Đây là năm vùng nội dung được FHSM chọn làm trọng tâm để triển khai luận và hiểu. Cách tổ chức này kế thừa nền tảng NCTM 2000, nhưng không lặp lại nguyên xi năm chuẩn nội dung của NCTM 2000. Thay vào đó, FHSM tái sắp xếp nội dung theo nhu cầu nhận thức ở bậc THPT: ghép số với đo lường để nhấn mạnh luận định lượng; tách đại số thành kí hiệu đại số và hàm số để làm rõ hai vùng luận khác nhau; giữ hình học như vùng đặc biệt của biện giải và chứng minh; và đặt thống kê-xác suất như vùng luận với dữ liệu, biến thiên và bất định.

    Điểm cần chú ý là trọng tâm của danh sách này không nằm ở việc phân loại lại chương trình thành năm phần cố định, mà ở việc chỉ ra những nơi luận và hiểu có thể được phát triển một cách rõ nét trong toán học THPT. Vì thế, các mục như ``Luận với số và đo lường'', ``Luận với kí hiệu đại số'', hay ``Luận với hàm số'' không phải là tên các chủ đề mới được thêm vào chương trình. Chúng là cách gọi cho những vùng nội dung quen thuộc khi được nhìn dưới lăng kính luận và hiểu: học sinh không chỉ học công thức, thủ tục hay kết quả, mà học cách nhận ra quan hệ, diễn giải ý nghĩa, kiểm soát phép biến đổi, xây dựng biện giải và kết nối kiến thức mới với tri thức đã có.

    Cách gọi ``Luận với kí hiệu đại số'' cần được hiểu như một lựa chọn có dụng ý. FHSM không nói chung là ``Luận với đại số'', vì trong đại số phổ thông, hoạt động luận của học sinh thường diễn ra thông qua việc đọc, viết, diễn giải và biến đổi kí hiệu: biến, biểu thức, phương trình, công thức, dấu bằng và các phép biến đổi đại số. Học sinh không tiếp cận đại số như một khối nội dung trừu tượng; các em tiếp cận nó qua hệ kí hiệu dùng để biểu thị quan hệ, khái quát hóa quan hệ và thao tác trên quan hệ ấy.

    Điểm này đặc biệt quan trọng vì kí hiệu đại số rất dễ bị dạy và học như chuỗi thao tác hình thức. Học sinh có thể biết chuyển vế, khai triển, rút gọn, phân tích thành nhân tử, giải phương trình, nhưng không thật sự hiểu vì sao các phép biến đổi ấy hợp lệ hoặc biểu thức đang biểu thị quan hệ nào. Vì thế, ``luận với kí hiệu đại số'' nhấn mạnh rằng việc học đại số không chỉ là thao tác đúng trên kí hiệu, mà là hiểu ý nghĩa của kí hiệu, kiểm soát phép biến đổi kí hiệu và dùng kí hiệu như công cụ để biện giải.}
  }

\TriPara
  {
    \EN{Within each content area, a number of key elements provide a broad structure for thinking about how the content area can be focused on reasoning and sense making. These key elements are not intended to be an exhaustive list of specific topics to be addressed. Instead, they provide a lens through which to view the potential of high school programs for promoting and developing mathematical reasoning and sense making. The task of creating a curriculum that fully achieves the goals of this publication will be challenging. Although important content must be addressed, this task requires much more than developing lists of topics to be taught in a particular course.}
  }
  {
    \VI{Trong mỗi mảng nội dung, một số yếu tố then chốt tạo nên cấu trúc rộng để suy nghĩ về cách đặt trọng tâm của mảng ấy vào luận và hiểu. Những yếu tố then chốt này không nhằm đưa ra danh sách đầy đủ các chủ đề cụ thể cần được giảng dạy. Đúng hơn, chúng cung cấp lăng kính để nhìn vào tiềm năng của chương trình trung học phổ thông trong việc thúc đẩy và phát triển luận và hiểu toán học. Việc xây dựng một chương trình thật sự đạt được các mục tiêu của ấn phẩm này sẽ là nhiệm vụ đầy thách thức. Dù nội dung quan trọng vẫn phải được xử lí, nhiệm vụ ấy đòi hỏi nhiều hơn rất nhiều so với việc lập ra danh sách chủ đề cần dạy trong một khóa học cụ thể.}
  }
  {
    \NOTE{Ở đây, \emph{key elements} không phải là danh sách đầy đủ các chủ đề cụ thể cần dạy, mà là những yếu tố giữ vai trò then chốt trong việc định hướng cách nhìn và cách triển khai nội dung theo hướng thúc đẩy luận và hiểu. ``Then chốt'' ở đây không nên hiểu theo nghĩa cấu thành bản chất, mà theo nghĩa tổ chức và định hướng: đây là những điểm tựa giúp người viết chương trình và người dạy nhận ra trong một mảng nội dung thì nên đặt trọng tâm vào đâu, nên nhìn tiềm năng của mảng ấy từ góc độ nào, và nên chọn những loại ví dụ, nhiệm vụ, hay mối liên hệ nào để luận và hiểu thật sự hiện diện trong việc học toán.}
  }

\TriPara
  {
    \EN{Furthermore, those responsible for making decisions about the curriculum must be aware of the need to offer experiences with reasoning and sense making within a broad curriculum that meets the needs of a wide range of students, preparing them for future success as citizens and in the workplace, as well as for careers in mathematics and science. Hard choices will be need to made about the topics traditionally included in the curriculum to make room for areas that have often been underemphasized, such as statistics. Many examples exist of schools and teachers who have already begun to reexamine their curricular priorities and instructional emphases to maximize focus on reasoning and sense making for all students across the years of high school mathematics. All schools and teachers must begin, or continue, this journey.}
  }
  {
    \VI{Hơn nữa, những người chịu trách nhiệm đưa ra quyết định về chương trình học phải ý thức được nhu cầu cung cấp trải nghiệm với luận và hiểu trong một chương trình rộng, đáp ứng nhu cầu của nhiều đối tượng học sinh khác nhau, chuẩn bị để các em thành công trong tương lai với tư cách công dân, trong môi trường làm việc, cũng như trong các nghề nghiệp thuộc toán học và khoa học. Sẽ cần đưa ra những lựa chọn khó khăn đối với các chủ đề vốn được đưa vào chương trình theo truyền thống, nhằm tạo chỗ cho những mảng thường bị xem nhẹ, chẳng hạn như thống kê. Đã có nhiều ví dụ về các trường và giáo viên bắt đầu xem xét lại ưu tiên chương trình và trọng tâm dạy học của mình để đặt trọng tâm tối đa vào luận và hiểu cho mọi học sinh trong suốt những năm học toán trung học phổ thông. Tất cả trường học và giáo viên phải bắt đầu, hoặc tiếp tục, hành trình này.}
  }{\NOTE{}}

\TriPara
  {
    \EN{Discrete mathematics, an active branch of contemporary mathematics that is widely used in business and industry, is an additional area of mathematics that is addressed in this publication. We argue that discrete mathematics should be an integral part of the school mathematics curriculum (NCTM 2000a). As in Principles and Standards, the main topics of discrete mathematics are distributed throughout the strands rather than receive separate treatment, as shown in figure 3.1.}
  }
  {
    \VI{Toán rời rạc, một nhánh toán học đương đại đang phát triển mạnh và được sử dụng rộng rãi trong kinh doanh và công nghiệp, là một mảng toán học bổ sung được đề cập trong ấn phẩm này. Chúng tôi cho rằng toán rời rạc nên là phần không thể thiếu của chương trình toán học phổ thông (NCTM 2000a). Cũng như trong \emph{Principles and Standards}, các chủ đề chính của toán rời rạc được phân bố xuyên suốt các mạch nội dung thay vì được xử lí như chủ đề tách biệt, như minh họa trong Hình 3.1.}
  }{\NOTE{}}

\TriFigure
  [
    Cần phân biệt hai hướng dịch gần nhau của họ từ \emph{recursion}: ``đệ quy'' và ``truy hồi''. Trong tiếng Anh, \emph{recursion} là khái niệm rộng, chỉ một cách xác định đối tượng, hàm, thủ tục hoặc quá trình bằng cách quay lại cùng kiểu đối tượng, cùng hàm hoặc cùng thủ tục ở cấp đơn giản hơn. Khi trọng tâm nằm ở cấu trúc tự quy chiếu hoặc tự gọi lại, tiếng Việt nên dùng ``đệ quy''. Vì vậy, các cụm như \emph{recursive function}, \emph{recursive algorithm}, \emph{recursive procedure}, \emph{recursive call}, \emph{recursively defined function} nên dịch lần lượt là ``hàm đệ quy'', ``thuật toán đệ quy'', ``thủ tục đệ quy'', ``lời gọi đệ quy'', ``hàm được định nghĩa đệ quy''.

    Tuy nhiên, trong toán học, họ từ \emph{recursion} cũng thường xuất hiện khi nói về dãy số, mẫu hình, quy luật lặp hoặc quá trình rời rạc, nơi giá trị ở bước sau được xác định từ một hoặc nhiều giá trị ở các bước trước. Khi ấy, tiếng Việt nên dùng ``truy hồi'', vì trọng tâm không phải là hành vi tự gọi lại của một hàm hay thủ tục, mà là quan hệ quay về các giá trị trước để xác định giá trị sau. Chẳng hạn, dãy
    \[
    u_1=1,\qquad u_{n+1}=2u_n+3
    \]
    nên gọi là ``dãy truy hồi'', vì số hạng $u_{n+1}$ được xác định từ số hạng $u_n$ theo một quy tắc. Theo hướng này, \emph{recursive sequence} nên dịch là ``dãy truy hồi'', \emph{recursive formula} là ``công thức truy hồi'', \emph{recursive rule} là ``quy tắc truy hồi'', \emph{recursively defined sequence} là ``dãy được xác định truy hồi''.

    Sự phân biệt này càng rõ hơn khi xét từ \emph{recurrence}. Trong tiếng Anh, \emph{recurrence} thường chỉ sự xuất hiện lặp lại hoặc quan hệ lặp lại; trong toán học, cụm \emph{recurrence relation} chỉ hệ thức xác định một số hạng, trạng thái hoặc giá trị rời rạc thông qua các số hạng, trạng thái hoặc giá trị trước đó. Vì vậy, \emph{recurrence relation} nên dịch ổn định là ``hệ thức truy hồi'', không phải ``hệ thức đệ quy''. Nói ngắn gọn: ``đệ quy'' làm nổi rõ cấu trúc tự quy chiếu hoặc tự gọi lại; ``truy hồi'' làm nổi rõ quan hệ bước sau--bước trước trong một dãy, mẫu hình hoặc quá trình rời rạc.

    Ở đây, \emph{recursion} nên dịch là ``truy hồi'' vì nó được đặt trong yếu tố then chốt ``nhiều biểu diễn của hàm số''. Trọng tâm không phải là thuật toán hay hàm tự gọi lại chính nó, mà là cách biểu diễn một hàm số, dãy số hoặc quá trình rời rạc bằng quy tắc xác định bước sau từ bước trước. Vì vậy, ``truy hồi'' làm rõ hơn ý nghĩa sư phạm của mục này: học sinh cần thấy rằng bên cạnh lời văn, bảng giá trị, đồ thị hoặc công thức tường minh, một quan hệ hàm số cũng có thể được biểu diễn bằng quy tắc truy hồi.
  ]
  {../figures/fhsm_2009_p19_f1.pdf}
  {../figures/fhsm_2009_p19_f1_vi.pdf}
  {Discrete mathematics in Focus in High School Mathematics}
  {{Toán rời rạc trong \emph{Trọng tâm trong toán học trung học phổ thông}}}
  {fig: fig_3.1}

\TriPara
  {
    \EN{The following chapters describe how reasoning and sense making can be promoted within the five content areas outlined above and characterize the key elements in each content area. These chapters include a series of examples in a variety of formats-including classroom vignettes and assessment items, as well as mathematical exposition. The examples are intended to provide an idealized illustration of how reasoning and sense making might unfold in high school mathematics.}
  }
  {
    \VI{Các chương tiếp theo mô tả cách luận và hiểu có thể được thúc đẩy trong năm mảng nội dung nêu trên, đồng thời mô tả các yếu tố then chốt trong mỗi mảng. Các chương này bao gồm một loạt ví dụ ở nhiều định dạng, trong đó có tình huống lớp học, câu hỏi đánh giá, cũng như diễn giải toán học. Những ví dụ này nhằm cung cấp một minh họa được lí tưởng hóa về cách luận và hiểu có thể diễn ra trong toán học trung học phổ thông.}
  }{\NOTE{}}

\CloseGrid